\documentclass[11pt,a4paper]{article}
\usepackage[utf8]{vietnam}
\usepackage[margin=2cm]{geometry}
\usepackage{amsmath,amssymb}
\usepackage{enumitem}
\usepackage{xcolor}
\usepackage{listings}
\usepackage{hyperref}
\usepackage{titlesec}
\usepackage{booktabs}
\usepackage{array}
\newcolumntype{L}[1]{>{\raggedright\arraybackslash}p{#1}}
\newcolumntype{C}[1]{>{\centering\arraybackslash}p{#1}}
\renewcommand{\arraystretch}{1.32}
\usepackage[htt]{hyphenat}
\emergencystretch=1.5em
\usepackage{setspace}
\setstretch{1.22}
\setlength{\parskip}{0.45em}
\definecolor{primarycolor}{RGB}{0, 80, 180}
\hypersetup{
    colorlinks=true,
    linkcolor=blue!70!black,
    urlcolor=blue!70!black
}

\titleformat{\section}{\large\bfseries\color{primarycolor}}{}{0em}{}[\vspace{0.18em}{\color{primarycolor}\titlerule}]
\titleformat{\subsection}{\normalsize\bfseries\color{primarycolor}}{}{0em}{}
\titlespacing*{\section}{0pt}{1.15em}{0.45em}
\titlespacing*{\subsection}{0pt}{0.75em}{0.25em}

\setlist[itemize]{itemsep=0.28em, topsep=0.28em, leftmargin=1.5em}
\setlist[enumerate]{itemsep=0.32em, topsep=0.32em, leftmargin=1.5em}

\lstdefinestyle{customc}{
    backgroundcolor=\color{gray!8},
    basicstyle=\small\ttfamily,
    breaklines=true,
    keywordstyle=\color{blue!80!black}\bfseries,
    commentstyle=\color{green!50!black}\itshape,
    frame=single,
    rulecolor=\color{gray!35},
    numbers=none,
    xleftmargin=0.8em,
    framexleftmargin=0.6em,
    showstringspaces=false
}

\begin{document}

\begin{center}
    {\Large \textbf{Đồ Án CV04: Building A Chinese-Vietnamese Parallel Corpus On Vietnamese Literature (Poetry)}}\\[0.3em]
    {\large \textbf{Môn học:} Xử lý ngôn ngữ tự nhiên (NLP)}\\[0.4em]
    \textbf{Nhóm sinh viên thực hiện:}\\[0.3em]
    \begin{tabular}{ll}
        1. Lê Gia Khánh & MSSV: 24122005 \\
        2. Trần Trung Hiếu & MSSV: 24122033 \\
        3. Phạm Chấn Khoa & MSSV: 24122038 \\
        4. Cái Hoàng Bảo Kin & MSSV: 24122040 \\
    \end{tabular}
\end{center}

\vspace{0.25em}
\noindent\textbf{Tóm tắt kết quả:}
Báo cáo tổng hợp đầy đủ \textbf{Mục 4} (\textit{Số lượng cặp, số dòng Hán -- Việt và kết quả tự kiểm tra}), \textbf{Mục 5} (\textit{Thông tin nguồn tài liệu và công cụ xử lý}) và \textbf{Mục 9} (\textit{5 tiêu chí đề xuất phương pháp}) theo tài liệu Hướng dẫn. Kho ngữ liệu hợp nhất của nhóm \texttt{CV04} đạt \textbf{8.122 cặp hợp lệ} (vượt chỉ tiêu $\ge 8.000$ cặp), tỷ lệ cảnh báo QA \textbf{6,43\%} ($\le 30\%$), \textbf{0 lỗi cấu trúc/thiếu vế} (\texttt{MISSING = 0}, \texttt{FORMAT\_ERROR = 0}) và đồng bộ \textbf{100\%} giữa \texttt{.tsv} với \texttt{.xlsx}.


% ============================================================================
\section{1. Tổng Quan Đồ Án \& Phân Công Thực Hiện Nhóm CV04}
% ============================================================================

Nhóm \texttt{CV04} xây dựng kho ngữ liệu song ngữ Hán -- Việt cấp câu/đoạn từ \textbf{03 bộ sách PDF scan} (\texttt{VPO\_01}, \texttt{VPO\_02}, \texttt{VPO\_03}; tổng \textbf{1.996 trang PDF gốc}). Với nguồn dữ liệu scan cần OCR, chỉ tiêu bắt buộc của nhóm 4 thành viên là tối thiểu \textbf{2.000 cặp/người} ($\ge \mathbf{8.000}$ \textbf{cặp hợp lệ} toàn nhóm).

\subsection{1.1. Quy chuẩn thư mục bàn giao và định dạng 4 cột}

Mỗi tác phẩm được đóng gói trong một thư mục \texttt{CV04\_VPO\_01/}, \texttt{CV04\_VPO\_02/}, \texttt{CV04\_VPO\_03/} gồm 04 thư mục con chuẩn hóa:

\begin{itemize}
    \item \textbf{\texttt{01\_nguon/}:} Chứa \textbf{1.584 ảnh PNG (300 DPI)} đặt tên theo số trang PDF 4 chữ số (\texttt{0007.png}, \dots) kèm tệp văn bản bổ sung trang 25 của \texttt{VPO\_03}.
    \item \textbf{\texttt{02\_ocr\_tho/}:} Lưu kết quả OCR thô tách riêng \texttt{han/*.json} (\textbf{1.584 tệp}) và \texttt{viet/*.json} (\textbf{1.584 tệp}) kèm tọa độ \textit{bounding box} và độ tin cậy.
    \item \textbf{\texttt{03\_text\_sach/}:} Gồm \texttt{han.txt} và \texttt{viet.txt} (\texttt{UTF-8 NFC}) đã làm sạch, gộp chung các bài của từng cuốn, phân tách bằng dòng tiêu đề và đồng bộ 1--1 theo dòng với tệp dóng hàng.
    \item \textbf{\texttt{04\_dong\_hang/}:} Gồm \texttt{<MA\_TAC\_PHAM>\_parallel.tsv} (phân tách bằng 1 ký tự \texttt{TAB}, không header), \texttt{<MA\_TAC\_PHAM>\_parallel.xlsx} (sheet \texttt{DONG\_HANG} khớp 100\% với \texttt{.tsv}, kèm sheet \texttt{THONG\_KE}), \texttt{qa\_report.json} và \texttt{word\_coverage\_report.json}.
\end{itemize}

\noindent\textbf{Quy tắc dóng hàng 4 cột:}
\begin{itemize}
    \item Cấu trúc cố định 4 cột: \textbf{\texttt{Hán}} -- \textbf{\texttt{Dịch âm}} -- \textbf{\texttt{Dịch nghĩa}} -- \textbf{\texttt{Dịch thơ}}.
    \item Mỗi bài thơ/từ/văn xuôi có \textbf{01 dòng tiêu đề riêng} ngay trước thân bài, đặt đúng cột tương ứng.
    \item Cột khuyết ở nguyên bản để trống \texttt{""} (bài thiếu dịch nghĩa ở \texttt{VPO\_01}; văn xuôi không có dịch thơ ở \texttt{VPO\_03}); phần \textit{Dịch theo nguyên điệu} của \texttt{VPO\_02} đặt vào cột \texttt{Dịch thơ}.
    \item Chữ Hán mờ/khuyết được bảo toàn bằng ký hiệu \texttt{[MASK]}, không tự đoán chữ.
\end{itemize}

\subsection{1.2. Phân công khối lượng 4 thành viên}

Toàn bộ \textbf{1.584 trang hợp lệ} được chia đều cho 4 thành viên (393--398 trang/người), giữ nguyên cặp trang chẵn -- lẻ đối diện trong \texttt{VPO\_01}:

\vspace{0.15em}
\begin{table}[h!]
\centering
\small
\begin{tabular}{L{3.5cm} C{1.3cm} L{2.8cm} L{2.8cm} L{2.8cm} C{1.6cm}}
\toprule
\textbf{Thành viên} & \textbf{Nhánh} & \textbf{CV04\_VPO\_01} & \textbf{CV04\_VPO\_02} & \textbf{CV04\_VPO\_03} & \textbf{Tổng} \\
\midrule
Phạm Chấn Khoa      & \texttt{khoa}  & 95 tr (\texttt{PDF 7..101})   & 104 tr (\texttt{PDF 134..237}) & 199 tr (\texttt{PDF 17..216})  & \textbf{398 tr} \\
Trần Trung Hiếu     & \texttt{hieu}  & 94 tr (\texttt{PDF 102..195}) & 104 tr (\texttt{PDF 238..341}) & 199 tr (\texttt{PDF 217..415}) & \textbf{397 tr} \\
Lê Gia Khánh        & \texttt{khanh} & 94 tr (\texttt{PDF 196..289}) & 103 tr (\texttt{PDF 342..444}) & 199 tr (\texttt{PDF 416..647}) & \textbf{396 tr} \\
Cái Hoàng Bảo Kin   & \texttt{kin}   & 92 tr (\texttt{PDF 290..381}) & 103 tr (\texttt{PDF 445..547}) & 198 tr (\texttt{PDF 648..861}) & \textbf{393 tr} \\
\midrule
\textbf{Hợp nhất (\texttt{main})} & \texttt{main} & \textbf{375 trang} & \textbf{414 trang} \textit{(411 align)} & \textbf{795 trang} & \textbf{1.584 tr} \\
\bottomrule
\end{tabular}
\caption{Phân công trang OCR và xử lý dữ liệu cho 4 thành viên nhóm \texttt{CV04}.}
\end{table}


% ============================================================================
\section{2. Thông Tin Nguồn Tài Liệu \& Kiểm Kê Phạm Vi (Yêu Cầu Mục 5.1)}
% ============================================================================

Nhóm đã kiểm kê đầu -- giữa -- cuối cả 3 tệp PDF nguồn, xác định công thức lệch trang giữa trang sách (\texttt{book}) và trang PDF (\texttt{pdf}), đồng thời chuẩn hóa phạm vi xử lý trong \texttt{config/book\_manifest.json}:

\vspace{0.15em}
\begin{table}[h!]
\centering
\small
\begin{tabular}{L{2.7cm} L{4.0cm} L{4.2cm} L{4.2cm}}
\toprule
\textbf{Trường kiểm kê} & \textbf{CV04\_VPO\_01 (\texttt{VPO\_01})} & \textbf{CV04\_VPO\_02 (\texttt{VPO\_02})} & \textbf{CV04\_VPO\_03 (\texttt{VPO\_03})} \\
\midrule
\textbf{Tên tác phẩm} & \textit{Hán Văn Tinh Túy} & \textit{Thể loại từ Việt Nam thời trung đại} & \textit{Thơ văn Lý -- Trần (Tập II -- Quyển Thượng)} \\
\addlinespace[0.2em]
\textbf{Tác giả / Chủ biên} & Lãng Nhân (Phùng Tất Đắc) & Phạm Văn Ánh / Nguyễn Xuân Diện -- Nguyễn Quảng Tuân & Nguyễn Huệ Chi (Chủ biên) \\
\addlinespace[0.2em]
\textbf{Tệp PDF nguồn} & \texttt{Hán văn tinh tuý - Lãng Nhân.pdf} & \texttt{Thể loại từ VN thời trung đại - Phạm Văn Ánh.pdf} & \texttt{Thơ văn Lý-Trần tập 2 - Nguyễn Huệ Chi.pdf} \\
\addlinespace[0.2em]
\textbf{Định dạng \& Số trang} & PDF Scan -- 386 trang & PDF Scan (lỗi text layer chữ Hán) -- 645 trang & PDF Scan -- 965 trang \\
\addlinespace[0.2em]
\textbf{Độ lệch trang} & $\texttt{pdf} = \texttt{book} + 6$ & $\texttt{pdf} = \texttt{book} + 7$ \textit{(sách 1..593)} & $\texttt{book} = \texttt{pdf} + 2$ \textit{(từ PDF 7)} \\
\addlinespace[0.2em]
\textbf{Phạm vi hợp lệ} & Sách tr.\,1..375 (\texttt{PDF 7..381}): \textbf{375 trang} & Phần Ba, sách tr.\,127..540 (\texttt{PDF 134..547}): \textbf{414 tr OCR} (\textbf{411 tr} dóng hàng) & Sách tr.\,19..863 (\texttt{PDF 17..861}): \textbf{795 trang} \\
\addlinespace[0.2em]
\textbf{Bố cục song ngữ} & Trang chẵn: \textit{Hán -- Phiên âm -- Dịch nghĩa}; trang lẻ đối diện: \textit{Dịch thơ} / tiêu đề. & Tuần tự theo khối: \textit{Hán -- Phiên âm -- Dịch nghĩa -- Dịch theo nguyên điệu -- Chú thích}. & Thơ 2 cột (\textit{Hán} trái, \textit{Phiên âm} phải) hoặc văn xuôi tuần tự -- \textit{Dịch nghĩa -- Dịch thơ -- Khảo đính}. \\
\bottomrule
\end{tabular}
\caption{Bảng kiểm kê 3 nguồn tài liệu song ngữ Hán -- Việt của nhóm \texttt{CV04}.}
\end{table}

\subsection{2.1. Các vùng loại trừ và xử lý ngoại lệ}

\begin{enumerate}
    \item \textbf{\texttt{CV04\_VPO\_01} (\textit{Hán Văn Tinh Túy}):}
    \begin{itemize}
        \item Loại trừ bìa, lời nói đầu (\texttt{PDF 1..6}) và mục lục cuối sách (\texttt{PDF 382..386}).
        \item Bài \textit{Thạch Hào lại} (\texttt{PDF 100}) và \textit{Thu dạ hoài ngâm} (\texttt{PDF 367..381}) nguyên bản không có phần dịch nghĩa xuôi: giữ đầy đủ cặp câu và để trống cột \texttt{Dịch nghĩa}.
    \end{itemize}

    \item \textbf{\texttt{CV04\_VPO\_02} (\textit{Thể loại từ Việt Nam thời trung đại}):}
    \begin{itemize}
        \item Chỉ lấy \textbf{Phần Ba} (\texttt{PDF 134..547}); bỏ bìa/mục lục (\texttt{PDF 1..7}), Phần 1--2 nghiên cứu thuần Việt (\texttt{PDF 8..133}) và Phần Phụ lục (\texttt{PDF 548..645}, gồm từ phổ và 118 bài ở Phụ lục 3 chỉ có chữ Hán chưa dịch).
        \item Loại trừ 3 trang \texttt{PDF 378..380} (04 bài từ chữ Nôm của Phạm Thái chỉ có phiên âm Quốc ngữ, khuyết chữ gốc và dịch nghĩa).
        \item Với 03 bài tàn khuyết một phần chữ Hán (\textit{Tây giang nguyệt} ở \texttt{PDF 138..140}, \textit{Quế chi hương} ở \texttt{PDF 221..225}, \textit{Hành hương tử} ở \texttt{PDF 303..304}), nhóm dóng hàng phần còn nguyên và giữ ký hiệu \texttt{[MASK]} / \texttt{[...]}.
    \end{itemize}

    \item \textbf{\texttt{CV04\_VPO\_03} (\textit{Thơ văn Lý -- Trần Tập II}):}
    \begin{itemize}
        \item Loại trừ phần đầu sách (\texttt{PDF 1..16}), phần khảo đính/mục lục cuối sách (\texttt{PDF 874..965}) và \textbf{04 tác phẩm Chữ Nôm}: \textit{Cư trần lạc đạo phú} (\texttt{PDF 493..518}), \textit{Đắc thú lâm tuyền thành đạo ca} (\texttt{PDF 519..525}), \textit{Vịnh Vân-yên tự phú} (\texttt{PDF 688..703}), \textit{Giáo tử phú} (\texttt{PDF 862..873}).
        \item \textbf{Bù trang trắng \texttt{PDF 25} (sách tr.\,27):} Tích hợp tệp \texttt{bai\_03\_thien\_tong\_chi\_nam\_tu\_dich\_nghia.txt} do chị Lưu Như cung cấp vào \texttt{01\_nguon/} để khôi phục trọn vẹn phần dịch nghĩa của bài 3 (\textit{Thiền tông chỉ nam tự}).
    \end{itemize}
\end{enumerate}


% ============================================================================
\section{3. Công Cụ Xử Lý \& Kiến Trúc Pipeline Tự Động (Yêu Cầu Mục 5.2)}
% ============================================================================

Hệ thống gồm \textbf{14 module Python} trong \texttt{src/}, điều phối qua CLI \texttt{main.py} và giao diện hiệu đính \texttt{web\_app/app.py}:

\vspace{0.15em}
\begin{table}[h!]
\centering
\small
\begin{tabular}{L{3.1cm} L{4.2cm} L{7.8cm}}
\toprule
\textbf{Công đoạn} & \textbf{Công cụ \& Module chính} & \textbf{Chức năng kỹ thuật} \\
\midrule
\textbf{1. Trích xuất \& Tiền xử lý ảnh} &
\texttt{PyMuPDF}, \texttt{OpenCV}, \texttt{Pillow}\newline (\texttt{src/pdf\_extractor.py}) &
Trích xuất ảnh PNG \textbf{300 DPI}; cắt \textbf{7\% lề trên} (\textit{running header}, số trang); dùng hình thái học phát hiện đường kẻ ngang nửa dưới trang để cắt bỏ chú thích chân trang (\texttt{VPO\_02}). \\
\addlinespace[0.25em]
\textbf{2. Nhận dạng quang học \& Khử chồng lấp} &
\texttt{PaddleOCR (chinese\_cht)} + \texttt{VietOCR (vgg\_transformer)}\newline (\texttt{src/ocr\_engine.py}) &
Chạy trên GPU Kaggle; tự động đối soát tọa độ giữa hai luồng Hán -- Việt (\texttt{reconcile\_han\_viet\_items}) để loại bỏ hộp OCR chồng lấp và gộp các cụm chữ Hán bị cắt ngang trên cùng dòng. \\
\addlinespace[0.25em]
\textbf{3. Làm sạch \& Khử nhầm nét Hán} &
\texttt{unicodedata (NFC)}, \texttt{Regex}, M.E.D. Levenshtein\newline (\texttt{src/text\_cleaner.py}, \texttt{src/ocr\_disambiguator.py}, \texttt{src/phonetic\_verifier.py}) &
Xóa chỉ mục chú thích \texttt{[1]}, \texttt{(1)}, dừng trước \texttt{KHẢO ĐÍNH / CHÚ THÍCH / LỜI BÀN}; đối chiếu Phiên âm bằng phép giao $S_1 \cap S_2$ trên \textbf{105 cụm nhầm nét} và \textbf{764 âm Hán -- Việt lõi} để sửa chữ Hán sai; phóng chiếu dấu câu từ Phiên âm sang câu Hán. \\
\addlinespace[0.25em]
\textbf{4. Phân đoạn \& Dóng hàng 4 cột} &
Quy hoạch động Gale--Church + Neo từ vựng/số đếm\newline (\texttt{src/work\_segmenter.py}, \texttt{src/aligner.py}, \texttt{src/corpus\_pipeline.py}) &
Tách bài xuyên trang kèm dòng tiêu đề riêng; tự động khôi phục phân vùng \texttt{Dịch nghĩa} khi OCR khuyết nhãn tiêu đề; tách mệnh đề song hành văn xuôi/biền ngẫu và dóng hàng 4 cột (\texttt{1--1}, \texttt{1--n}, \texttt{n--1}, \texttt{n--m}). \\
\addlinespace[0.25em]
\textbf{5. Xuất tệp, Kiểm định QA \& Hiệu đính} &
\texttt{src/exporter\_and\_qa.py}, \texttt{src/word\_aligner.py}, \texttt{web\_app/app.py} &
Xuất đồng thời \texttt{.tsv} và \texttt{.xlsx} (khớp 100\%); kiểm tra 6 nhóm lỗi QA, tính độ phủ Hán tự; hỗ trợ rà soát đối chiếu ảnh gốc trên giao diện web Streamlit. \\
\bottomrule
\end{tabular}
\caption{Ngăn xếp công cụ và kiến trúc pipeline xử lý của nhóm \texttt{CV04}.}
\end{table}


% ============================================================================
\section{4. Báo Cáo Số Lượng Cặp, Số Dòng Hán -- Việt \& Kết Quả Tự Kiểm Tra (Yêu Cầu Mục 4)}
% ============================================================================

\subsection{4.1. Thống kê số lượng trang, số dòng Hán, số dòng Việt và số cặp dóng hàng}

\vspace{0.15em}
\begin{table}[h!]
\centering
\small
\begin{tabular}{L{5.6cm} C{2.1cm} C{2.1cm} C{2.1cm} C{2.5cm}}
\toprule
\textbf{Chỉ số thống kê} & \textbf{CV04\_VPO\_01} & \textbf{CV04\_VPO\_02} & \textbf{CV04\_VPO\_03} & \textbf{Tổng toàn nhóm} \\
\midrule
Tổng số trang PDF gốc & 386 & 645 & 965 & \textbf{1.996 trang} \\
Số trang ảnh PNG (\texttt{01\_nguon/}) & 375 & 414 & 795 & \textbf{1.584 trang} \\
Số tệp JSON OCR Hán / Việt (\texttt{02\_ocr\_tho/}) & 375 / 375 & 414 / 414 & 795 / 795 & \textbf{1.584 / 1.584} \\
Số trang hợp lệ đưa vào dóng hàng & 375 & 411 & 795 & \textbf{1.581 trang} \\
\midrule
Số hộp dòng OCR thô luồng Hán (\texttt{han/*.json}) & 10.505 & 13.635 & 27.561 & \textbf{51.701 dòng} \\
Số hộp dòng OCR thô luồng Việt (\texttt{viet/*.json}) & 5.458 & 8.205 & 24.781 & \textbf{38.444 dòng} \\
\midrule
\textbf{Số dòng văn bản Hán sạch (\texttt{03\_text\_sach/han.txt})} & \textbf{1.265} & \textbf{2.059} & \textbf{4.798} & \textbf{8.122 dòng} \\
\textbf{Số dòng văn bản Việt sạch (\texttt{03\_text\_sach/viet.txt})} & \textbf{1.265} & \textbf{2.059} & \textbf{4.798} & \textbf{8.122 dòng} \\
\midrule
\textbf{Tổng số cặp dóng hàng hợp lệ (\texttt{.tsv} \& \texttt{.xlsx})} & \textbf{1.265} & \textbf{2.059} & \textbf{4.798} & \textbf{8.122 cặp} \\
\quad -- \textit{Số dòng tiêu đề bài} & 25 & 174 & 285 & \textbf{484 dòng} \\
\quad -- \textit{Số dòng nội dung thân bài} & 1.240 & 1.885 & 4.513 & \textbf{7.638 dòng} \\
\quad -- \textit{Số dòng có bản Dịch thơ (Cột 4)} & 1.124 & 1.321 & 1.472 & \textbf{3.917 dòng} \\
\bottomrule
\end{tabular}
\caption{Thống kê số lượng trang, số dòng Hán, số dòng Việt và số cặp dóng hàng hợp lệ.}
\end{table}

\begin{itemize}
    \item \textbf{Vượt chỉ tiêu sản lượng:} Toàn nhóm đạt \textbf{8.122 cặp hợp lệ} (trung bình \textbf{2.030,5 cặp/thành viên}, vượt ngưỡng $\ge 8.000$ cặp).
    \item \textbf{Đồng bộ tuyệt đối:} \texttt{03\_text\_sach/han.txt} (\textbf{8.122 dòng}) và \texttt{viet.txt} (\textbf{8.122 dòng}) khớp 1--1 theo từng dòng vật lý với \texttt{<MA\_TAC\_PHAM>\_parallel.tsv} và sheet \texttt{DONG\_HANG} của \texttt{.xlsx}.
\end{itemize}

\subsection{4.2. Kết quả tự kiểm tra chất lượng (QA Validation \& Word Coverage)}

Module \texttt{src/exporter\_and\_qa.py} kiểm định toàn bộ \textbf{8.122 cặp câu} theo công thức Mục 8 Hướng dẫn:
\[
\text{Tỷ lệ lỗi} = \frac{\text{Số cặp có cảnh báo lỗi}}{\text{Tổng số cặp được kiểm tra hợp lệ}} \quad (\text{Ngưỡng nghiệm thu: } \le 30\%; \text{ Mục tiêu nội bộ: } \le 20\%)
\]

\vspace{0.15em}
\begin{table}[h!]
\centering
\small
\begin{tabular}{L{5.6cm} C{2.1cm} C{2.1cm} C{2.1cm} C{2.5cm}}
\toprule
\textbf{Chỉ tiêu tự kiểm tra (QA \& Coverage)} & \textbf{CV04\_VPO\_01} & \textbf{CV04\_VPO\_02} & \textbf{CV04\_VPO\_03} & \textbf{Tổng / Trung bình} \\
\midrule
Tổng số cặp kiểm tra hợp lệ (\texttt{valid\_pairs}) & 1.265 & 2.059 & 4.798 & \textbf{8.122 cặp} \\
Số cặp sạch hoàn toàn (\texttt{clean\_pairs}) & 1.225 (96,84\%) & 1.964 (95,39\%) & 4.411 (91,93\%) & \textbf{7.600 (93,57\%)} \\
Số cặp có cảnh báo QA (\texttt{error\_pairs}) & 40 & 95 & 387 & \textbf{522 cặp} \\
\textbf{Tỷ lệ cảnh báo lỗi QA (\texttt{error\_rate})} & \textbf{3,16\%} & \textbf{4,61\%} & \textbf{8,07\%} & \textbf{6,43\% ($\le 30\%$)} \\
\midrule
\quad 1. Lỗi \texttt{MISSING} (thiếu vế Hán hoặc Việt) & \textbf{0} & \textbf{0} & \textbf{0} & \textbf{0} \\
\quad 2. Lỗi \texttt{FORMAT\_ERROR} (lỗi cấu trúc Tab/dòng) & \textbf{0} & \textbf{0} & \textbf{0} & \textbf{0} \\
\quad 3. Lỗi \texttt{LEFTOVER\_MARKER} (sót chỉ mục chú thích) & \textbf{0} & \textbf{0} & \textbf{0} & \textbf{0} \\
\quad 4. Cảnh báo \texttt{OMISSION} (tỷ lệ dài Việt/Hán $< 0,7$) & 33 & 66 & 158 & \textbf{257 lượt} \\
\quad 5. Cảnh báo \texttt{ADDITION} (tỷ lệ dài Việt/Hán $> 4,5$) & 6 & 16 & 133 & \textbf{155 lượt} \\
\quad 6. Cảnh báo \texttt{NUMBER\_MISMATCH} (lệch từ chỉ số) & 1 & 17 & 108 & \textbf{126 lượt} \\
\midrule
Độ phủ từ vựng Hán tự trung bình (\texttt{han\_coverage}) & 59,49\% & 47,22\% & 43,30\% & \textbf{46,82\%} \\
Điểm tin cậy dóng hàng trung bình (\texttt{avg\_confidence}) & 0,5744 & 0,5037 & 0,4816 & \textbf{0,5017} \\
Kiểm tra đồng bộ nội dung \texttt{TSV == XLSX} & Khớp 100\% & Khớp 100\% & Khớp 100\% & \textbf{Khớp 100\%} \\
Kết luận nghiệm thu tự động (\texttt{is\_valid}) & \textbf{ĐẠT} & \textbf{ĐẠT} & \textbf{ĐẠT} & \textbf{ĐẠT (6,43\%)} \\
\bottomrule
\end{tabular}
\caption{Kết quả tự kiểm tra chất lượng (QA) và độ phủ từ vựng trên 8.122 cặp câu.}
\end{table}

\noindent\textbf{Nhận xét kết quả QA:}
\begin{enumerate}
    \item \textbf{Không có lỗi cấu trúc và thiếu vế:} Cả 3 bộ sách đạt \texttt{MISSING = 0}, \texttt{FORMAT\_ERROR = 0}, \texttt{LEFTOVER\_MARKER = 0}; mỗi dòng \texttt{.tsv} có đúng 4 cột phân tách bởi 3 ký tự \texttt{TAB} và khớp 100\% với sheet \texttt{DONG\_HANG} của \texttt{.xlsx}.
    \item \textbf{Bản chất các cảnh báo mềm (6,43\%):} 522 dòng cảnh báo chủ yếu đến từ đặc thù dịch văn học cổ điển:
    \begin{itemize}
        \item \texttt{OMISSION} (257 lượt): Câu Hán chứa nhiều hư từ (\textit{chi, hồ, giả, dã, nhi, tắc}) được dịch cô đọng nên tỷ lệ số từ Việt / số chữ Hán $< 0,7$.
        \item \texttt{ADDITION} (155 lượt): Câu Hán ngắn (2--4 chữ chứa điển cố/thuật ngữ Thiền học) được dịch giả diễn giải chi tiết hoặc ghép chung nhiều câu dịch thơ (\texttt{1--n}), làm tỷ lệ dài $> 4,5$.
        \item \texttt{NUMBER\_MISMATCH} (126 lượt): Chữ số mang nghĩa ước lệ trong điển cố (\textit{cửu thiên, vạn nhất, bách lộc, tứ phương}) được dịch thoát ý sang tiếng Việt.
    \end{itemize}
\end{enumerate}


% ============================================================================
\section{5. Thuyết Minh Đề Xuất Phương Pháp Khác Theo 5 Tiêu Chí (Mục 9 Hướng Dẫn)}
% ============================================================================

Theo Mục 9 của Hướng dẫn, nhóm đề xuất quy trình tự động hóa kết hợp hậu kiểm gồm 4 cải tiến: (i) \textbf{Tiền xử lý cắt chú thích bằng hình thái học \& đối soát khử hộp OCR chồng lấp Hán -- Việt}, (ii) \textbf{Khử nhầm nét chữ Hán và phóng chiếu dấu câu qua cầu nối Phiên âm Hán -- Việt}, (iii) \textbf{Dóng hàng 4 cột kết hợp tách mệnh đề song hành và quy hoạch động có neo từ vựng}, và (iv) \textbf{Giao diện hiệu đính CAT Workbench}.

\subsection{5.1. Tiêu chí 1: Lý do lựa chọn phương pháp}

\begin{itemize}
    \item \textbf{Xử lý triệt để lỗi mất text layer, nhiễu trang scan và hộp OCR chồng lấp:} Do PDF \texttt{VPO\_02} lỗi font mất chữ Hán và cả 3 sách có nhiều tiêu đề đầu trang/chú thích chân trang, việc kết xuất ảnh 300 DPI kèm cắt 7\% lề trên, tách đường kẻ ngang chân trang và đối soát tọa độ hộp Hán -- Việt (\texttt{reconcile\_han\_viet\_items}) giúp loại bỏ hoàn toàn nhiễu nhận dạng chéo giữa hai luồng OCR.
    \item \textbf{Tận dụng ``cầu nối Phiên âm'' để sửa lỗi OCR chữ Hán:} OCR chữ Hán cổ dễ nhầm các chữ đồng hình dị tự, trong khi dòng \textit{Phiên âm Quốc ngữ} đi kèm được \texttt{VietOCR} nhận dạng rất chính xác. Module \texttt{src/ocr\_disambiguator.py} dùng khoảng cách chỉnh sửa tối thiểu (\textit{Weighted Levenshtein}) khớp từng âm tiết phiên âm với chữ Hán trên phép giao $S_1 \cap S_2$ (\textbf{105 cụm nhầm nét}, \textbf{764 âm Hán -- Việt lõi}) để tự động sửa chữ Hán nhận dạng sai.
    \item \textbf{Dóng hàng 4 cột chính xác cho cả thơ và văn xuôi cổ:} Các công cụ dóng hàng 2 cột sẵn có không xử lý được cấu trúc 4 cột và dễ trượt dòng ở văn biền ngẫu câu dài. Thuật toán tách mệnh đề song hành (\texttt{\_align\_prose\_unit}) kết hợp quy hoạch động Gale--Church có điểm thưởng neo từ vựng Hán -- Việt (\texttt{src/word\_aligner.py}) giúp khớp đồng thời cả 4 cột kèm dòng tiêu đề từng bài.
\end{itemize}

\subsection{5.2. Tiêu chí 2: Dữ liệu đầu vào và đầu ra}

\begin{itemize}
    \item \textbf{Đầu vào (Input):} 03 tệp PDF scan gốc trong \texttt{raw\_data/} (1.996 trang), tệp văn bản bù trang 25 của \texttt{VPO\_03}, cùng 02 tệp cấu hình \texttt{config/book\_manifest.json} và \texttt{config/member\_splits.json}.
    \item \textbf{Đầu ra (Output):} 03 thư mục bàn giao \texttt{CV04\_VPO\_01..03/} gồm 1.584 ảnh PNG 300 DPI (\texttt{01\_nguon/}, 964,5 MB), 1.584 cặp JSON OCR thô (\texttt{02\_ocr\_tho/}), 03 cặp \texttt{han.txt} và \texttt{viet.txt} (\texttt{03\_text\_sach/}, 8.122 dòng/bên), và 03 cặp \texttt{<MA\_TAC\_PHAM>\_parallel.\{tsv,xlsx\}} (\texttt{04\_dong\_hang/}, \textbf{8.122 cặp hợp lệ}) kèm báo cáo QA.
\end{itemize}

\subsection{5.3. Tiêu chí 3: Cách kiểm chứng chất lượng}

\begin{enumerate}
    \item \textbf{Kiểm định QA tự động (\texttt{python main.py validate}):} Rà soát 6 nhóm lỗi trên toàn bộ 8.122 cặp, xác nhận \texttt{MISSING = 0}, \texttt{FORMAT\_ERROR = 0}, tỷ lệ cảnh báo \textbf{6,43\%} ($\le 30\%$).
    \item \textbf{Đối chiếu đồng bộ \texttt{TSV == XLSX}:} Tự động đọc ngược tệp \texttt{.tsv} và sheet \texttt{DONG\_HANG} của \texttt{.xlsx} sau khi xuất để kiểm chứng khớp 100\% từng ô.
    \item \textbf{Đo độ phủ từ vựng cấp từ (\texttt{src/word\_aligner.py}):} Ghi nhận độ phủ Hán tự trung bình đạt \textbf{59,49\%} (\texttt{VPO\_01}), \textbf{47,22\%} (\texttt{VPO\_02}) và \textbf{43,30\%} (\texttt{VPO\_03}).
    \item \textbf{Kiểm thử tự động (\texttt{tests/test\_suite.py}):} 08 kịch bản unit test \& integration test cô lập trong thư mục tạm, kiểm tra toàn bộ pipeline (\textbf{8/8 bài test đạt 100\%}).
    \item \textbf{Hậu kiểm trực quan (\texttt{web\_app/app.py}):} Lọc các dòng cảnh báo trên Streamlit CAT Workbench để đối chiếu ảnh trang sách gốc và chỉnh sửa trực tiếp.
\end{enumerate}

\subsection{5.4. Tiêu chí 4: Rủi ro và phương án dự phòng}

\vspace{0.15em}
\begin{table}[h!]
\centering
\small
\begin{tabular}{L{3.8cm} L{4.8cm} L{6.5cm}}
\toprule
\textbf{Rủi ro kỹ thuật} & \textbf{Nguyên nhân trong nguồn} & \textbf{Phương án dự phòng đã áp dụng} \\
\midrule
\textbf{1. Nhiễu chú thích \& khảo đính tràn vào nội dung} &
\texttt{VPO\_02} nhiều chú thích chân trang; \texttt{VPO\_03} có mục \textit{Khảo đính, Dị bản, Chú thích} sau mỗi bài. &
Cắt chân trang theo đường kẻ ngang bằng OpenCV (\texttt{VPO\_02}); dùng regex không dấu ngắt tự động khi gặp tiêu đề \texttt{KHẢO ĐÍNH / CHÚ THÍCH / LỜI BÀN} (\texttt{VPO\_03}). \\
\addlinespace[0.25em]
\textbf{2. Trang scan trắng hoặc bài khuyết bản gốc/bản dịch} &
\texttt{VPO\_03} trắng trang \texttt{PDF 25}; \texttt{VPO\_02} có 4 bài khuyết chữ gốc (\texttt{PDF 378..380}) và 3 bài khuyết một phần; \texttt{VPO\_01} có bài thiếu dịch nghĩa. &
Tích hợp tệp bù trang 25 của giảng viên; loại trừ trang khuyết hoàn toàn trong \texttt{book\_manifest.json}; để trống cột khuyết \texttt{""} và giữ \texttt{[MASK]} cho chữ mờ. \\
\addlinespace[0.25em]
\textbf{3. Khác biệt môi trường (Windows/macOS) \& thiếu GPU local} &
4 thành viên dùng hệ điều hành khác nhau; máy cá nhân không cài sẵn CUDA hoặc \texttt{openpyxl}. &
Chạy OCR trên Kaggle GPU và lưu JSON thô vào Git (\texttt{02\_ocr\_tho/}); các khâu làm sạch, dóng hàng, xuất \texttt{.xlsx} (OpenXML tự viết) và QA chạy thuần Python chuẩn không cần GPU. \\
\bottomrule
\end{tabular}
\caption{Rủi ro kỹ thuật và phương án dự phòng trong pipeline của nhóm \texttt{CV04}.}
\end{table}

\subsection{5.5. Tiêu chí 5: Khả năng tái lập bởi người khác (Reproducibility)}

Toàn bộ kết quả dóng hàng và kiểm định QA có thể tái lập 100\% trên mọi máy tính cài Python 3 qua các lệnh CLI chuẩn:

\begin{lstlisting}[style=customc]
# 1. Cai dat thu vien phu thuoc (PyMuPDF, Pillow, openpyxl, streamlit)
pip install -r requirements.txt

# 2. (Tuy chon) Trich xuat anh PNG 300 DPI tu 3 file PDF goc
python main.py extract --book ALL --pages valid --dpi 300

# 3. Dong hang 4 cot toan bo 3 bo sach tu JSON OCR tho (02_ocr_tho/)
python main.py align --book ALL

# 4. Kiem dinh chat luong QA, do phu tu vung va doi chieu TSV == XLSX
python main.py validate

# 5. Chay 8 bai kiem thu tu dong va mo giao dien hieu dinh CAT Workbench
python -m unittest discover -s tests -v
python main.py web
\end{lstlisting}

\end{document}
